% Generated draft: requires mathematical and source review.
\documentclass[11pt]{article}
\usepackage[margin=25mm]{geometry}
\usepackage{fontspec}
\setmainfont{Latin Modern Roman}
\usepackage{amsmath,amssymb,mathtools,unicode-math}
\setmathfont{Latin Modern Math}
\usepackage{enumitem,needspace,xurl,hyperref}
\setlength{\parindent}{0pt}
\setlength{\parskip}{5pt}
\setlength{\emergencystretch}{3em}
\newcommand{\N}{\mathbb{N}}
\newcommand{\Z}{\mathbb{Z}}
\newcommand{\Q}{\mathbb{Q}}
\newcommand{\R}{\mathbb{R}}
\newcommand{\C}{\mathbb{C}}
\newcommand{\F}{\mathbb{F}}
\newcommand{\A}{\mathbb{A}}
\newcommand{\PP}{\mathbb{P}}
\newcommand{\E}{\mathbb{E}}
\newcommand{\eps}{\varepsilon}
\newcommand{\dd}{\,\mathrm{d}}
\begin{document}
\section*{UnsolvedMath ID: 30003038}
\textbf{Source problem code:} catalogue entry
\subsection{Definitions and mathematical statement}
The supplied question concerns comparison of shapes, including cardiac configurations at different times and different hand postures, and the choice of a scale-invariant distance along curves in shape space. It is a construction and modelling question; the record does not specify enough structure to determine a unique mathematical proposition.

Write \(\mathcal S\) for a prospective space of shapes. A family of configurations would be represented by a map \(\gamma:T\to\mathcal S\), where the parameter domain \(T\), the representation of configurations, and any regularity requirements remain unspecified. In particular, the record does not determine whether the objects being compared are points \(\gamma(t)\) along such a curve or entire curves \(\gamma\).

If ``distance'' is intended to mean a metric on a specified comparison space \(X\), its definition would require a function \(d:X\times X\to[0,\infty)\) satisfying, for all \(x,y,z\in X\),
\[
 d(x,y)=0\iff x=y,\qquad d(x,y)=d(y,x),\qquad d(x,z)\leq d(x,y)+d(y,z).
\]
Here \(X\) could be a shape space, a quotient space, or a space of curves; the source does not choose among these possibilities or require a genuine metric rather than another dissimilarity measure.

If positive rescaling is represented by maps \(D_\lambda:X\to X\), \(\lambda>0\), with \(D_1=\operatorname{id}_X\) and \(D_\lambda D_\mu=D_{\lambda\mu}\), two possible meanings of scale invariance are
\[
 d(D_\lambda x,D_\lambda y)=d(x,y)
 \quad(x,y\in X,\ \lambda>0)
\]
and the stronger requirement
\[
 d(D_\lambda x,D_\mu y)=d(x,y)
 \quad(x,y\in X,\ \lambda,\mu>0).
\]
These are possible interpretations, not additional hypotheses supplied by the record. The latter requirement is incompatible with a genuine metric on a representation space containing distinct rescalings of the same object; identifying such rescalings or allowing a pseudometric would require an explicit modelling choice.

\textbf{Source question.} How should the indicated configurations be represented and compared, and what scale-invariant distance should be used along the resulting curves in shape space? A precise construction problem requires specification of the objects, admissible curves, equivalence relations, meaning of scale invariance, and criteria for a suitable comparison. The record supplies none of these choices sufficiently to select a single formal target. It also leaves unresolved whether ``along'' requests pointwise distances, path lengths, an induced intrinsic distance, or a distance between entire curves.
\subsection{Short English statement}
The question asks how to compare changing configurations, such as cardiac shapes or hand postures, using a scale-invariant distance in shape space. The representation of shapes, treatment of scaling, and meaning of distance along a curve must be specified before this becomes a definite mathematical problem.
\subsection{Sources}
\begin{itemize}
\item Local UnsolvedMath registry: supplied complete record with permanent ID \(30003038\), source problem code \texttt{catalogue entry}, in the Oberwolfach Reports Open Problems collection. The record gives proposed year 2016 and supplies no named proposer.
\item Source report: \emph{Mathematical Imaging and Surface Processing} (2016), Oberwolfach Reports, pp.\ 155--214, DOI 10.4171/owr/2016/4. Supplied report link and literature reference labelled ``Oberwolfach Report 4/2016'': \url{https://doi.org/10.4171/owr/2016/4}. No specific passage within these pages is identified in the supplied record.
\item Supplied literature reference: Srivastava, Klassen, Joshi, and Jermyn, \emph{Shape analysis of elastic curves}: \url{https://pubmed.ncbi.nlm.nih.gov/20921581/}.
\item Supplied literature assessment, dated exactly \texttt{2026-08-21}: the record labels the problem partially solved and states that elastic shape analysis provides practical Riemannian distances for curves with invariance under reparameterization and standard shape transformations, while no unique universal scale-invariant distance answers this broad natural-shape question. This is supplied dated curation derived from the OWR record, not an independently verified literature assessment; its status and cited claims require independent verification.
\item The registry labels the statement \texttt{corrected\_verified} and reports a manual semantic, typographic, and LaTeX review. These are supplied provenance claims, not independent verification of the formulation or mathematical status of this draft.
\end{itemize}
\end{document}
